In the Bucegi mountains, part of the Carpathian mountains, after the end of
the First World War an iron cross was built by the former King of Romania
called Ferdinand the I-st and his wife Queen Maria. The cross is an unique
monument in Europe. The monument is an impressive iron cross called ``The
Heros' Cross'' which in 2013 entered in the Guiness Book as the cross build on
the highest altitude mountain peek.

The cross was built on the plane plateau situated on the top of the peek
called Caraiman, at the altitude $H = 2300$\,m from sea level. Its height,
including the base-support is $h = 39.3$\,m. The horizontal arms of the cross
are oriented on the N--S direction. The latitude of the Cross is
$\varphi = 45^\circ$.

\begin{description}
\item[A.] In the evening of 21st of March 2014, the summer equinox day, two
  eagles stop from their flight, first near the monument, and the second, on
  the top of the Cross as seen in figure 1. The two eagles are on the same
  vertical direction. The sky was very clear, so the eagles could see the
  horizon and observe the Sun set. Each eagle start to fly right in the moment
  that each of it observes that the Sun completely disappears.

  In the same time, an astronomer located at the sea-level, at the base of the
  Bucegi Mountains. Assume that he is on the same vertical with the two
  eagles.

  \ioaafig{2014_TH_L1-f1.pdf}

  Assuming negligible the atmospheric refraction, solve the following
  questions:
  \begin{description}
  \item[1)] Calculate the duration of the sunset, measured by the astronomer.
  \item[2)] Calculate the durations of sunset measured by each of the two
    eagles and indicate which of the eagles leaves first the Cross. What is
    the time interval between the leaving moments of the two eagles.
  \end{description}
  \emph{The following information is necessary:} The duration of the sunset
  measurement starts when the solar disc is tangent to the horizon line and
  stops when the solar disc completely disappears. The Earth's rotation period
  is $T_E = 24$\,h, the radius of the Sun $R_S = 6.96 \times 10^{5}$\,km,
  Earth--Sun distance $d_{ES} = 15 \times 10^{7}$\,km, the local latitude of
  the Heroes Cross $\varphi = 45^\circ$.
\item[B)] At a certain moment of the next day, 22nd March 2014, the two eagles
  come back to the Heroes Cross. One of the eagles lands on the top of the
  vertical pillar of the Cross and the other one land on the horizontal
  plateau, just in the end point of the shadow of the vertical pillar of the
  Cross.
  \begin{description}
  \item[1)] Calculate the distance between the two eagles, if this distance
    has the minimum possible value.
  \item[2)] Calculate the length of the horizontal arms of the Cross $l_b$, if
    the shadow on the plateau of one of the arm of the cross has the length
    $u_b = 7$\,m.
  \end{description}
\item[C)] At midnight, the astronomer visit the cross and, from the top of it,
  he identifies a bright star at the limit of the circumpolarity. He named
  this star ``Eagles Star''. Knowing that due to the atmospheric refraction
  the horizon lowering is $\xi = 34'$, calculate:
  \begin{description}
  \item[1)] The ``Eagles star'' declination;
  \item[2)] The ``Eagles star'' maximum height above the horizon.
  \end{description}
\end{description}
